\subsection{预序关系}

设 \(\mathbb{R}\{x, y\}\) 是一个关系，其中 \(x\) 与 \(y\) 是不同的字母。如果 \(\mathbb{R}\) 是传递的，并且如果我们有 \(\mathbb{R}\{x, y\} \Longrightarrow (\mathbb{R}\{x, x\} \text{ 且 } \mathbb{R}\{y, y\})\) ，并不必然推出 \(\mathbb{R}\) 是一个序关系，因为关系

\[
(\mathbf {R} \{x, y \} \text { 且 } \mathbf {R} \{y, x \})
\]

并不必然蕴含 \(x = y\) . \(\mathbb{R}\{x, y\}\) 被称为 \(x\) 与 \(y\) 之间的预序关系； \(\mathbb{R}\{y, x\}\) 则也是 \(x\) 与 \(y\) 之间的预序关系，称为关系 \(\mathbb{R}\{x, y\}\) 的相反关系。

例如，设 R 为 P(E) 的覆盖 E 的子集所构成的集合（第二章，§4，第 6 号）。R 的元素 R、R' 之间的关系“R 比 R' 更粗”（第二章，§4，第 6 号）是传递且自反的。但两个不同的覆盖可能使得每一个都比另一个更粗；例如，当 R' 是（在 P(E) 中）R 与 E 的一个子集的并集，且该子集包含于 R 的某个集合中但不属于 R 时，就是这种情况。

但在任何情况下，关系 \((\mathbb{R}\{x, y\}\) 并且 \(\mathbb{R}\{y, x\})\) 是一个等价关系 \(\mathbb{S}\{x, y\}\) （关于 \(x\) 与 \(y\) ）. 设 \(x', y'\) 为不同于 \(x, y\) 且不出现在 \(\mathbb{R}\) 中的字母。于是 \(\mathbb{R}\{x, y\}\) （关于 \(x\) 与 \(y\) ）与等价关系 \(\mathbb{S}\{x, x'\}\) 及 \(\mathbb{S}\{y, y'\}\) 相容；换言之（第二章，§6，第8号），关系

\[
(\mathrm{R} \{x, y \} \text {   且   } \mathrm{S} \{x, x ^ {\prime} \} \text {   且   } \mathrm{S} \{y, y ^ {\prime} \})
\]

蕴含R \(\{x', y'\}\) .

集合 E 上的一个预序关系是一个预序关系 R\{x, y\} 使得关系 R\{x, x\} 等价于 \(x \in E\) 。关系 R\{x, y\} 于是蕴含“ \(x \in E\) 并且 \(y \in E\) ''.

如果 \(\mathbf{R}\{x, y\}\) 是集合 \(\mathbf{E}\) 上的一个预序关系，则上文定义的关系 \(\mathbf{S}\{x, y\}\) 是 \(\mathbf{E}\) 上的一个等价关系. 设 \(\mathbf{R}'\{\mathbf{X}, \mathbf{Y}\}\) 表示该关系

\[
\mathbf {X} \in \mathbf {E} / \mathrm{S} \text { 且 } \mathbf {Y} \in \mathbf {E} / \mathrm{S} \text { 且 } (\exists x) (\exists y) (x \in \mathbf {X} \text { 且 } y \in \mathbf {Y} \text { 且 } \mathbf {R} \{x, y \}),
\]

也就是说，R在过渡到商集（关于x和y）时所诱导的关系；我们在第二章§6第3号中看到，它与该关系等价。

\[
\mathrm{X} \in \mathrm{E} / \mathrm{S} \text {   且   } \mathrm{Y} \in \mathrm{E} / \mathrm{S} \text {   且   } (\forall x) (\forall y) ((x \in \mathrm{X} \text {   且   } y \in \mathrm{Y}) \implies \mathrm{R} \{x, y \}).
\]

让我们证明 \(R'\{X, Y\}\) 是E/S的元素之间的一个序关系。关系 \((R'\{X, Y\}\) 并且 \(R'\{Y, Z\})\) 等价于

\(\mathbf{X} \in \mathbf{E} / \mathbf{S}\) 并且 \(\mathbf{Y} \in \mathbf{E} / \mathbf{S}\) 并且 \(\mathbf{Z} \in \mathbf{E} / \mathbf{S}\) 并且 \((\forall x)(\forall y)(\forall z)((x \in \mathbf{X}\) 并且 \(y \in \mathbf{Y}\) 并且 \(z \in \mathbf{Z}) \Longrightarrow (\mathbf{R}\{x, y\}\) 并且 \(\mathbf{R}\{y, z\}))\)

（第一章，§4，准则C40，C41）。由于R\{x, y\} 是传递的，且Y ∈ E/S ⟶ Y ≠ φ（第二章，§6，第2号），由此立即得出R' \{X, Y\} 是传递的。其次，(R' \{X, Y\} 和R' \{Y, X\})等价于

\[
\begin{array}{l} \mathrm{X} \in \mathrm{E} / \mathrm{S} \text { 且 } \mathrm{Y} \in \mathrm{E} / \mathrm{S} \text { 且 } \\ (\forall x) (\forall y) ((x \in \mathrm{X} \text { 且 } y \in \mathrm{Y}) \Longrightarrow (\mathrm{R} \{x, y \} \text { 且 } \mathrm{R} \{y, x \})), \end{array}
\]

，因此等价于

\(X \in E/S\) 并且 \(Y \in E/S\) 并且 \((\forall x)(\forall y)((x \in X \text{ 且 } y \in Y) \Longrightarrow S\{x, y\})\) ，从而蕴含

\[
\mathrm{X} \in \mathrm{E} / \mathrm{S} \quad \text { 且 } \quad \mathrm{Y} \in \mathrm{E} / \mathrm{S} \quad \text { 且 } \quad \mathrm{X} = \mathrm{Y}.
\]

此外，R \(\{x, y\}\) 蕴含R \(\{x, x\}\) 和 R \(\{y, y\}\) ，因此R' \(\{X, Y\}\) 蕴含以下每个关系

\[
\begin{array}{l} \mathrm{X} \in \mathrm{E} / \mathrm{S} \text { 且 } (\forall x) ((x \in \mathrm{X}) \Longrightarrow \mathrm{R} \{x, x \}), \\ \mathrm{Y} \in \mathrm{E} / \mathrm{S} \text { 且 } (\forall y) ((y \in \mathrm{Y}) \Longrightarrow \mathrm{R} \{y, y \}), \end{array}
\]

由此 \(R'\{X,Y\}\) 蕴含 \((R'\{X,X\}\) 并且 \(R'\{Y,Y\})\) 。最后，由于 \(x \in E\) 蕴含 \(R\{x,x\}\) , \(X \in E/S\) 蕴含 \(R'\{X,X\}\) ，我们断言的证明至此完成。 \(R'\{X,Y\}\) 被称为与 \(R\{x,y\}\) 相伴的序关系。

集合E上的预序是一个对应关系 \(\Gamma = (G, E, E)\) 以E为源，以E为目标，并且满足 \((x, y) \in G\) 是 E 上的一个预序关系。为简便起见，我们有时将图 \(G\) （ \(\Gamma\) 的图）称为 E 上的一个预序。为此，必要且充分条件是 \(\Delta \subset G\) 并且 \(G \circ G \subset G\) （这意味着 \(G \circ G = G\) ）。等价关系 S

对应于预序关系 \((x, y) \in G\) 于是以 \(G \cap \overline{G}\) 为其图；与 \((x, y) \in G\) 相关联的序关系以子集 \(G'\) （ \((E/S) \times (E/S)\) 的子集）为图，该子集（第二章，§6，第8号）是 \(G\) 在典范映射下的像所对应的子集，该典范映射从 \(E \times E\) 到

\[
(\mathbf {E} \times \mathbf {E}) / (\mathbf {S} \times \mathbf {S}).
\]

示例。* 设A为具有单位元的环。关系 \((\exists z)(z \in A\) 并且 \(y = zx)\) （A 的两个元素 \(x\) , \(y\) 之间）是 A 上的一个预序关系；它读作“\(x\) 是 \(y\) 的右因子”或“\(y\) 是 \(x\) 的左倍数”。 *

